% =====================================================================
\section{Related work}\label{sec:related}
% =====================================================================
This section reviews four lines of work that NITSCI builds on: deep learning for DR grading, lesion-aware and
attention-based models, ordinal and imbalance-aware learning, and training strategies for small medical datasets.
Comprehensive surveys of deep learning for DR are available in \citet{tsiknakis2021deep}.

\subsection{Deep learning for DR detection and grading}
Early applications of convolutional neural networks (CNNs) to fundus photographs showed that end-to-end learning
could replace hand-crafted lesion detectors \citep{pratt2016convolutional,abramoff2016improved}. The study of
\citet{gulshan2016development} trained an Inception-style network on more than 100\,000 graded images and reported
specialist-level performance for referable DR, and similar results were obtained on independent and multi-ethnic
cohorts \citep{gargeya2017automated,ting2017development}. Later systems extended the task to the full disease spectrum
and to lesion-level outputs \citep{dai2021deep}. In parallel, general-purpose backbones improved steadily: residual
networks \citep{he2016deep}, compound-scaled EfficientNets \citep{tan2019efficientnet}, the faster-training
EfficientNetV2 family \citep{tan2021efficientnetv2}, modernised CNNs \citep{liu2022convnet}, and Vision Transformers
\citep{dosovitskiy2021image,liu2021swin}. Most recently, retinal foundation models pretrained by self-supervision on
hundreds of thousands of unlabelled images have improved data efficiency across retinal tasks
\citep{zhou2023foundation}.

Public benchmarks have made methods comparable. Messidor \citep{decenciere2014messidor}, DDR \citep{li2019ddr},
APTOS-2019 \citep{aptos2019} and IDRiD \citep{porwal2018idrid} differ in population, camera, image quality and
labelling protocol. IDRiD is particularly valuable because it combines image-level grades with pixel-level lesion
annotations for a subset of images, and its grading challenge established an official partition and leaderboard
\citep{porwal2020idrid}. Its small size, however, makes it a demanding test of data efficiency. Uncertainty-aware
grading \citep{araujo2020drgraduate} and studies of reference-standard quality \citep{krause2018grader} have further
shown that reported accuracy depends strongly on how labels and test sets are constructed.

\subsection{Lesion-aware and attention-based grading}
Because DR grades are defined by lesions, many methods try to make the network focus on them. Zoom-in-Net
\citep{wang2017zoom} produces attention maps that highlight suspicious regions and re-examines them at higher
resolution. Deep image mining \citep{quellec2017deep} derives lesion heat-maps from a grading network trained with
image-level labels only. Collaborative learning \citep{zhou2019collaborative} trains lesion segmentation and grading
jointly, using a small set of lesion masks to guide the grading network. CANet \citep{li2020canet} uses attention to
model the relation between DR and diabetic macular oedema, and CABNet \citep{he2021cabnet} introduces a category
attention block that learns discriminative regions for each grade under class imbalance. Lesion-aware transformers
\citep{sun2021lesion} learn lesion queries with a transformer decoder and use them to predict the grade.

Channel and spatial attention modules such as squeeze-and-excitation \citep{hu2018squeeze} and CBAM
\citep{woo2018cbam}, as well as non-local blocks \citep{wang2018nonlocal}, are generic ways to re-weight features.
Self-attention in Transformers \citep{vaswani2017attention} relates all positions at once, and cross-attention
between two token streams has been used to fuse multi-scale branches in CrossViT \citep{chen2021crossvit}. Dedicated
lesion segmentation networks, from U-Net \citep{ronneberger2015unet} to recent heterogeneous Mamba-based ensembles
\citep{wu2026helnet,gu2023mamba}, produce explicit lesion maps but require pixel-level annotations that are expensive
to obtain.

NITSCI sits between these families. Like deep image mining it learns a lesion-attention map from grade labels alone;
like CrossViT it uses cross-attention between two token streams; and like mixture-of-experts models
\citep{jacobs1991adaptive,shazeer2017outrageously} it combines the streams with a learned gate. The difference from
single-map attention modules is that the lesion map defines a \emph{second expert} whose tokens interact with the
global expert in both directions, rather than only re-weighting a single pooled feature.

\subsection{Ordinal and imbalance-aware learning}
Ordinal classification can be reduced to a series of binary ``greater than $k$'' problems
\citep{frank2001simple}, an idea later implemented with multiple CNN outputs \citep{niu2016ordinal} and made
rank-consistent by CORAL \citep{cao2020rank}. Other work optimises the evaluation metric directly, for example with a
differentiable weighted-kappa loss \citep{delatorre2018weighted}. In DR competitions a simple and robust alternative
has been to regress the grade as a real number and to optimise the cut-points between grades on validation data;
NITSCI follows this practice but combines the regression output with the expected value of the classification head
so that both training signals contribute to the decision.

Class imbalance is usually handled by re-weighting or re-sampling. Focal loss \citep{lin2017focal} down-weights easy
examples, and class-balanced loss \citep{cui2019class} re-weights classes by their effective number of samples.
Label smoothing \citep{szegedy2016rethinking} reduces over-confidence, which is useful when a few labels may be noisy.
We use inverse-square-root class weights with label smoothing, a moderate choice that does not over-amplify the
rarest grades.

\subsection{Learning from small medical datasets}
Transfer learning from ImageNet \citep{deng2009imagenet} is the default initialisation for medical image
classification. Better ImageNet models tend to transfer better \citep{kornblith2019better}, and larger pretraining
corpora such as ImageNet-21k help further \citep{ridnik2021imagenet21k}, but the benefit for medical images depends on
the domain gap and model size \citep{raghu2019transfusion}. An intermediate fine-tuning stage on a larger in-domain
dataset is a natural way to narrow this gap; we use APTOS-2019 for this purpose.

Several inference- and optimisation-level techniques improve robustness at little cost. Ensembles of independently
trained networks improve accuracy and uncertainty estimates \citep{lakshminarayanan2017simple}; averaging weights
along the training trajectory, through exponential moving averages \citep{polyak1992acceleration} or stochastic
weight averaging \citep{izmailov2018averaging}, finds flatter solutions that generalise better; and test-time
augmentation (TTA) reduces prediction variance in medical imaging \citep{wang2019aleatoric}. Modern networks are
often poorly calibrated \citep{guo2017calibration}, which matters when a probability threshold defines a referral
decision.

Finally, evaluation practice is critical in small-data settings. Selecting models or thresholds with the same data
used to estimate performance biases the estimate upward \citep{varma2006bias}, and data leakage is a widespread
source of irreproducible results in machine-learning-based science \citep{kapoor2023leakage}. Reporting guidelines
such as TRIPOD+AI \citep{collins2024tripod} ask for a clear separation of development and evaluation data,
confidence intervals and transparent reporting of all modelling choices. Our protocol is designed to meet these
requirements.
\sectionend{related}
